\documentclass[UTF8,a4paper,11pt,openany]{ctexbook}
\usepackage{../common/statlearnbook}
\SLSetBookMetadata{第 4 册}{统计学习方法讲义 v2.6.0}{无监督学习与矩阵分解}
\begin{document}
\setcounter{page}{644}
\setcounter{chapter}{28}
\setcounter{figure}{0}
\pagestyle{headings}
\chaptermark{潜在语义分析与非负矩阵分解}
\begin{optimizationbox}
\textbf{优化解释。}截断 SVD 对无约束秩 $k$ 问题给出全局最优解。NMF 的可行集要求两个因子均非负，联合目标并非凸；乘法更新不给出全局最优证书。
\end{optimizationbox}
如\cref{fig:V4-C05-two-geometries}所示，两条路线从同一文本向量出发，却把“低维”解释为不同的可行几何。
% v2.3.1 figure UID: FIG-P504-01
% Proposed post-v2.3.1 polish for FIG-P504-01: protect key-label size and stroke hierarchy.
\tikzset{
  slfig-FIG-P504-01/.style={font=\fontsize{9.4pt}{11.2pt}\selectfont,
    every path/.append style={line cap=round,line join=round}},
  slfig-FIG-P504-01-axis/.style={draw=SLGray,line width=.52pt,-{Stealth[length=1.6mm,width=1mm]}},
  slfig-FIG-P504-01-basis/.style={-{Stealth[length=2.1mm,width=1.25mm]},line width=1pt,draw=SLBlue},
  slfig-FIG-P504-01-target/.style={-{Stealth[length=2.1mm,width=1.25mm]},line width=1.02pt,draw=SLGold},
  slfig-FIG-P504-01-residual/.style={draw=SLGray,densely dashed,line width=.58pt},
  slfig-FIG-P504-01-right-angle/.style={draw=SLBlue,line width=.68pt},
  slfig-FIG-P504-01-reconstruction/.style={-{Stealth[length=1.9mm,width=1.15mm]},draw=SLTeal,line width=.96pt},
  slfig-FIG-P504-01-title/.style={font=\bfseries\fontsize{10.2pt}{12.2pt}\selectfont},
  slfig-FIG-P504-01-note/.style={font=\fontsize{9.2pt}{11pt}\selectfont,align=center,text=SLGray},
  slfig-FIG-P504-01-summary/.style={font=\fontsize{9.2pt}{11pt}\selectfont,
    draw=SLRuleGray,line width=.58pt,rounded corners=2pt,fill=SLSoftGray,
    minimum width=116mm,minimum height=6.5mm,align=center},
}
\begin{figure}[H]
\centering
\begin{tikzpicture}[slfig-FIG-P504-01,>=Stealth,
  every node/.style={font=\fontsize{9.4pt}{11.2pt}\selectfont},
]
  % 两面板共享同一目标向量 x 与秩 K=2。
  \begin{scope}[shift={(-3.55,0)}]
    \node[slfig-FIG-P504-01-title,text=SLBlue] at (0,2.62) {LSA：正交子空间（$K=2$）};
    \fill[SLBlue!6] (-2.15,-1.55) rectangle (2.15,1.70);
    \draw[slfig-FIG-P504-01-axis] (-2.35,0)--(2.45,0);
    \draw[slfig-FIG-P504-01-axis] (0,-1.75)--(0,1.95);
    \draw[slfig-FIG-P504-01-basis] (0,0)--(1.65,0) node[below right] {$u_1$};
    \draw[slfig-FIG-P504-01-basis] (0,0)--(0,1.45) node[above left] {$u_2$};
    \draw[slfig-FIG-P504-01-right-angle] (.23,0)--(.23,.23)--(0,.23);
    \draw[slfig-FIG-P504-01-target] (0,0)--(-1.05,1.28) node[above left] {$x$};
    \draw[slfig-FIG-P504-01-residual] (-1.05,1.28)--(-1.05,0);
    \fill[SLGold] (-1.05,0) circle (1.85pt);
    \node[anchor=north,text=SLGold] at (-1.05,-.06) {$\widehat x=U_2U_2^{\mathsf T}x$};
    \node[slfig-FIG-P504-01-note] at (0,-2.05)
      {$y=U_2^{\mathsf T}x$，坐标可正可负};
  \end{scope}

  \begin{scope}[shift={(3.55,0)}]
    \node[slfig-FIG-P504-01-title,text=SLTeal] at (0,2.62) {NMF：非负锥（$K=2$）};
    \fill[SLTeal,opacity=.09] (0,0)--(2.10,.48)--(.72,1.95)--cycle;
    \draw[slfig-FIG-P504-01-axis] (-2.35,0)--(2.45,0);
    \draw[slfig-FIG-P504-01-axis] (0,-1.75)--(0,1.95);
    \draw[slfig-FIG-P504-01-basis,draw=SLTeal] (0,0)--(2.10,.48) node[right] {$w_1$};
    \draw[slfig-FIG-P504-01-basis,draw=SLTeal] (0,0)--(.72,1.95) node[above] {$w_2$};
    % 与左图完全相同的目标向量坐标。
    \draw[slfig-FIG-P504-01-target] (0,0)--(-1.05,1.28) node[above left] {$x$};
    \draw[slfig-FIG-P504-01-residual] (-1.05,1.28)--(.72,1.20);
    \draw[slfig-FIG-P504-01-reconstruction]
      (0,0)--(.72,1.20) node[right] {$Wh$};
    \fill[SLTeal] (.72,1.20) circle (1.85pt);
    \node[slfig-FIG-P504-01-note] at (0,-2.05)
      {$h_1,h_2\ge0$，只允许锥内组合};
  \end{scope}

  \node[slfig-FIG-P504-01-summary] at (0,-2.82)
    {同秩 $K=2$\quad/\quad 不同约束与目标：正交投影 vs. 非负重构};
\end{tikzpicture}
\caption{LSA的正交子空间表示与NMF的非负锥表示：二者维数相同，坐标约束和最优性不同}\label{fig:V4-C05-two-geometries}
\end{figure}

\end{document}
